\subsection{DEFINITION OF A REAL NUMBER BY MEANS OF ITS EXPANSION}

Conversely, suppose we are given an integer \(q_0\) and a sequence \((v_n)\) ( \(n \geqslant 1\) ) of integers such that \(0 \leqslant v_n \leqslant a_n - 1\) ; we ask whether there is a number \(x\) whose expansion (4) is such that \(p_0 = q_0\) and \(u_n = v_n\) for all \(n\) . If such a number exists it is unique, because it is equal to \(q_0 + \sum_{n=1}^{\infty} \frac{v_n}{d_n}\) .

For each integer \(m > 0\) we have (by the principle of comparison)

\[
\sum_ {n = m + 1} ^ {\infty} \frac {v _ {n}}{d _ {n}} \leqslant \sum_ {n = m + 1} ^ {\infty} \frac {a _ {n} - 1}{d _ {n}} = \sum_ {n = m + 1} ^ {\infty} \left(\frac {1}{d _ {n - 1}} - \frac {1}{d _ {n}}\right) = \frac {1}{d _ {m}}
\]

and the extreme left-hand and right-hand terms are equal only if \(v_{n} = a_{n} - 1\) for each \(n > m\) (§ 7, no. 1, Theorem 2). Hence the series whose general term is \(\frac{v_n}{d_n}\) is convergent; moreover, if \(x = q_0 + \sum_{n=1}^{\infty} \frac{v_n}{d_n}\) , we have

\[
s _ {m} = q _ {0} + \sum_ {n = 1} ^ {m} \frac {v _ {n}}{d _ {n}} \leqslant x \leqslant s _ {m} + \frac {1}{d _ {m}}
\]

and \(x = s_m + \frac{1}{d_m}\) only if \(v_n = a_n - 1\) for each \(n > m\) . Since \(s_m\) is a fraction with denominator \(d_m\) , the approximation \(r_m\) to \(x\) to within \(1/d_m\) by defect is equal to \(s_m\) or \(s_m + \frac{1}{d_m}\) ; and the latter alternative can occur only if \(v_{n} = a_{n} - 1\) for all \(n > m\) . Thus:

\begin{enumerate}
\def\labelenumi{(\roman{enumi})}
\item
  There is an infinity of values of \(n\) such that \(v_{n} < a_{n} - 1\) : then the series \(q_{0} + \sum_{n=1}^{\infty} \frac{v_{n}}{d_{n}}\) is identical with the expansion of its sum \(x\) .
\item
  There is an integer \(m \geqslant 0\) such that \(v_{n} = a_{n} - 1\) whenever \(n > m\) , and \(v_{m} < a_{m} - 1\) (if \(m > 0\) ); then the sum \(x\) of the series \(q_{0} + \sum_{n=1}^{\infty} \frac{v_{n}}{d_{n}}\) is equal to the rational number
\end{enumerate}

\[
q _ {0} + \sum_ {n = 1} ^ {m} \frac {v _ {n}}{d _ {n}} + \frac {1}{d _ {m}}\tag{5}
\]

which is of the form \(k/d_{m}\) (k an integer); the expansion of x is identical with the series (5), in which all the terms with indices \textgreater{} m are zero; such an expansion is said to be terminating, or to terminate. The series

\[
q _ {0} + \sum_ {n = 1} ^ {\infty} \frac {v _ {n}}{d _ {n}} = q _ {0} + \sum_ {n = 1} ^ {m} \frac {v _ {n}}{d _ {n}} + \sum_ {n = m + 1} ^ {\infty} \frac {a _ {n} - 1}{d _ {n}}\tag{6}
\]

is called the improper expansion of the number x.

Conversely, let x be a rational number which can be written in the form of a fraction with denominator \(d_{n}\) for some value of n.~Let m be the smallest integer such that x is of the form \(k/d_{m}\) (k an integer); we have \(r_{n}<x\) for n\textless m, and \(r_{m}=x\) , and therefore the expansion of x is of the form (5), and x has an improper expansion given by (6); moreover this improper expansion is unique.

A rational number, written in its irreducible form \(p / q\) , is equal to a fraction with denominator \(d_n\) if and only if \(q\) divides \(d_n\) (the number \(m\) will therefore be the smallest integer \(n\) such that \(q\) divides \(d_n\) ). It can happen that every rational number has this property (for a suitably chosen \(n\) ): this will be the case if and only if every integer \(> 0\) divides some \(d_n\) , e.g., if \(d_n = n!\) . If the \(d_n\) have this property, then a number is rational if and only if its expansion, relative to the sequence \((d_n)\) , terminates.

To summarize: to every sequence \(s\) whose initial term \(q_{0}\) is an arbitrary integer, and whose term \(v_{n}(n \geqslant 1)\) is such that \(0 \leqslant v_{n} \leqslant a_{n} - 1\) , corresponds a real number equal to \(q_{0} + \sum_{n=1}^{\infty} \frac{v_{n}}{d_{n}}\) ; if \(I_{n}\) denotes the interval \([0, a_{n} - 1]\) of \(N\) , we thus define a mapping \(\varphi\) of \(X = Z \times \prod_{n=1}^{\infty} I_{n}\) onto the real line \(R\) ; moreover the equation \(\varphi(s) = x\) , where \(x \in R\) is given, has one solution if \(x\) is not a fraction with denominator \(d_{n}\) (for some \(n\) ), and two solutions otherwise.

